\section{Capítulo 4 -- Roteamento e endpoints}

No Capítulo 3, a rota \texttt{/api/health} mostrou que basta um arquivo para criar um \textit{endpoint}. Neste capítulo você vai entender a regra que liga cada URL a uma função do código, ler tudo o que chega em uma requisição (parâmetros, cabeçalhos e corpo), montar respostas com o status certo e criar as primeiras rotas de verdade do TaskFlow: listar, criar, buscar e excluir tarefas. No final, vai conhecer o que o Next.js faz sozinho, como as respostas a \texttt{OPTIONS} e o status \texttt{405}, e as formas práticas de testar uma API durante o desenvolvimento.

\subsection{4.1 Do endereço ao código}

\textbf{Roteamento} é o mapeamento entre uma requisição (método e URL) e a função que vai atendê-la. Em muitos \textit{frameworks}, como o Express, as rotas são escritas em código, uma por linha: \texttt{app.get("/api/tasks", listar)}. No Next.js, o mapeamento é feito pelas \textbf{pastas}: cada pasta dentro de \texttt{app/} é um pedaço da URL, e um arquivo \texttt{route.ts} transforma aquela pasta em \textit{endpoint}.

\begin{center}
\begin{tabular}{@{}P{0.47\linewidth}P{0.43\linewidth}@{}}
\toprule
\rowcolor{headergray}
\tblheader{Arquivo} & \tblheader{URL} \\
\midrule
\texttt{app/route.ts} & \texttt{/} \\
\rowcolor{rowgray} \texttt{app/api/health/route.ts} & \texttt{/api/health} \\
\texttt{app/api/tasks/route.ts} & \texttt{/api/tasks} \\
\rowcolor{rowgray} \texttt{app/api/tasks/[id]/route.ts} & \texttt{/api/tasks/42}, \texttt{/api/tasks/abc}... \\
\texttt{app/api/auth/login/route.ts} & \texttt{/api/auth/login} \\
\bottomrule
\end{tabular}
\end{center}

Dentro do \texttt{route.ts}, cada função exportada com o nome de um método HTTP atende àquele método: \texttt{GET}, \texttt{POST}, \texttt{PUT}, \texttt{PATCH}, \texttt{DELETE}, \texttt{HEAD} e \texttt{OPTIONS}. O nome precisa estar em maiúsculas. Assim, o mapeamento completo de um \textit{endpoint} é ``pasta + nome da função'': \texttt{POST /api/tasks} é a função \texttt{POST} do arquivo \texttt{app/api/tasks/route.ts}. Esses arquivos se chamam \textbf{\textit{Route Handlers}}.

Uma pasta pode ter um \texttt{route.ts} ou uma página (\texttt{page.tsx}), mas não os dois: o Next.js não saberia se deve devolver HTML ou executar a rota.

\subsection{4.2 Lendo a requisição}

Toda função de rota recebe a requisição como primeiro parâmetro. O tipo \texttt{NextRequest} é o \texttt{Request} padrão da web, o mesmo do \texttt{fetch}, com alguns atalhos a mais. Ela traz tudo o que o cliente enviou:

\begin{center}
\begin{tabular}{@{}P{0.3\linewidth}P{0.6\linewidth}@{}}
\toprule
\rowcolor{headergray}
\tblheader{O quê} & \tblheader{Como ler} \\
\midrule
\textit{query string} & \texttt{request.nextUrl.searchParams.get("status")} \\
\rowcolor{rowgray} cabeçalho & \texttt{request.headers.get("authorization")} \\
corpo em JSON & \texttt{await request.json()} \\
\rowcolor{rowgray} corpo em texto & \texttt{await request.text()} \\
corpo de formulário & \texttt{await request.formData()} \\
\rowcolor{rowgray} parte variável da URL & \texttt{(await contexto.params).id} (Seção 4.6) \\
método e URL & \texttt{request.method} e \texttt{request.url} \\
\bottomrule
\end{tabular}
\end{center}

Três detalhes evitam surpresas. O \texttt{searchParams.get} devolve sempre um texto, ou \texttt{null} se o parâmetro não veio: \texttt{?page=2} chega como \texttt{"2"}, e não como o número 2. O nome dos cabeçalhos não diferencia maiúsculas de minúsculas, então \texttt{get("Authorization")} e \texttt{get("authorization")} dão no mesmo. E o corpo só pode ser lido \textbf{uma vez}: depois de um \texttt{request.json()}, um \texttt{request.text()} falha.

\subsection{4.3 Montando a resposta}

A função devolve um \texttt{Response}, também o padrão da web. As três formas que aparecem em quase toda API são estas:

\begin{codigo}
\begin{Verbatim}[breaklines=true, breaksymbolleft={}, fontsize=\footnotesize]
// 200 com corpo JSON (o status padrão é 200)
return Response.json({ data: tarefas });

// outro status e cabeçalhos extras
return Response.json(tarefa, {
  status: 201,
  headers: { Location: `/api/tasks/${tarefa.id}` },
});

// sem corpo: o 204 não pode ter corpo
return new Response(null, { status: 204 });
\end{Verbatim}
\end{codigo}

O \texttt{Response.json} transforma o objeto em texto JSON e acrescenta o cabeçalho \texttt{Content-Type: application/json}. Você vai encontrar também o \texttt{NextResponse.json}, de \texttt{next/server}, que faz o mesmo e tem recursos extras para \textit{cookies} e redirecionamentos; para uma API, o \texttt{Response} padrão basta.

\subsection{4.4 Um usuário provisório}

Cada pessoa só pode ver as próprias tarefas, então toda rota precisa saber quem está fazendo a requisição. Isso depende da autenticação, que chega no Capítulo 9. Até lá, uma função provisória finge que toda requisição vem da usuária de teste criada pelo \textit{seed}:

\begin{codigo}
\begin{Verbatim}[breaklines=true, breaksymbolleft={}, fontsize=\footnotesize]
// lib/usuario-atual.ts -- PROVISÓRIO: até o Capítulo 9, toda
// requisição age como a usuária de teste criada pelo seed.
import { db } from "@/lib/db";

export async function usuarioAtual() {
  return db.user.findUniqueOrThrow({
    where: { email: "demo@taskflow.com" },
    select: { id: true, name: true },
  });
}
\end{Verbatim}
\end{codigo}

O \texttt{findUniqueOrThrow} busca pelo e-mail, que é único, e lança um erro se não encontrar. O \texttt{select} traz só o identificador e o nome: a senha nunca precisa sair do banco aqui. As rotas vão chamar \texttt{usuarioAtual()} sem saber como ela funciona por dentro. Quando a autenticação estiver pronta, só esta função muda, e nenhuma rota precisa ser alterada. Isolar o que vai mudar atrás de uma função é uma das técnicas mais úteis de organização de código.

\subsection{4.5 A coleção de tarefas}

O arquivo \texttt{app/api/tasks/route.ts} atende à coleção: listar (\texttt{GET}) e criar (\texttt{POST}).

\begin{codigo}
\begin{Verbatim}[breaklines=true, breaksymbolleft={}, fontsize=\footnotesize]
// app/api/tasks/route.ts -- a coleção de tarefas
import type { NextRequest } from "next/server";
import { db } from "@/lib/db";
import { TaskStatus } from "@/app/generated/prisma/client";
import { usuarioAtual } from "@/lib/usuario-atual";

// GET /api/tasks e GET /api/tasks?status=PENDING
export async function GET(request: NextRequest) {
  const usuario = await usuarioAtual();
  const status = request.nextUrl.searchParams.get("status");

  if (status && !Object.hasOwn(TaskStatus, status)) {
    return Response.json(
      { erro: "status deve ser PENDING, IN_PROGRESS ou DONE" },
      { status: 400 },
    );
  }

  const tarefas = await db.task.findMany({
    where: {
      userId: usuario.id,
      ...(status && { status: status as TaskStatus }),
    },
    orderBy: { createdAt: "desc" },
  });

  return Response.json({ data: tarefas });
}
\end{Verbatim}
\end{codigo}

O Prisma gera, para cada \texttt{enum} do esquema, um objeto com os valores permitidos: \texttt{TaskStatus} é \texttt{\{ PENDING: "PENDING", IN\_PROGRESS: ..., DONE: ... \}}. O \texttt{Object.hasOwn(TaskStatus, status)} confere se o texto recebido é uma dessas chaves. Repare que não foi usado \texttt{status in TaskStatus}: o operador \texttt{in} também enxerga as propriedades herdadas de todo objeto, e \texttt{?status=toString} passaria pela verificação e quebraria a consulta.

O filtro usa uma construção comum em JavaScript: \texttt{...(status \&\& \{ status \})}. Se \texttt{status} for \texttt{null}, a expressão vale \texttt{null}, e espalhar \texttt{null} em um objeto não acrescenta nada; se houver um valor, o campo \texttt{status} entra no \texttt{where}. O \texttt{as TaskStatus} diz ao TypeScript que o texto já foi conferido. Por fim, \texttt{orderBy} lista as tarefas mais recentes primeiro, e o resultado vai dentro de \texttt{data}, o envelope combinado no Capítulo 2. A criação fica no mesmo arquivo:

\begin{codigo}
\begin{Verbatim}[breaklines=true, breaksymbolleft={}, fontsize=\footnotesize]
// POST /api/tasks  -- corpo: {"title": "..."}
export async function POST(request: NextRequest) {
  const usuario = await usuarioAtual();
  const corpo = await request.json().catch(() => null);

  if (typeof corpo?.title !== "string" || !corpo.title.trim()) {
    return Response.json(
      { erro: "O título é obrigatório." },
      { status: 400 },
    );
  }

  const tarefa = await db.task.create({
    data: { title: corpo.title.trim(), userId: usuario.id },
  });

  return Response.json(tarefa, {
    status: 201,
    headers: { Location: `/api/tasks/${tarefa.id}` },
  });
}
\end{Verbatim}
\end{codigo}

Se o corpo não for um JSON válido, \texttt{request.json()} lança um erro, e a API responderia \texttt{500}, como se a falha fosse do servidor. O \texttt{.catch(() => null)} transforma esse caso em \texttt{null}, que cai na verificação seguinte e recebe um \texttt{400}, o código certo para um erro do cliente. O \texttt{corpo?.title} usa o \textbf{encadeamento opcional}: se \texttt{corpo} for \texttt{null}, o resultado é \texttt{undefined} em vez de um erro. A verificação aceita só um texto que não fique vazio depois do \texttt{trim}, que tira os espaços das pontas.

Esta verificação ``à mão'' cobre só o título. Conferir todos os campos (prioridade, data, tamanho do texto) desse jeito seria longo e fácil de errar; o Capítulo 6 substitui tudo por um esquema do Zod, e o Capítulo 7 completa a criação com os demais campos.

\subsection{4.6 Segmentos dinâmicos: uma tarefa}

Para atender \texttt{/api/tasks/42}, \texttt{/api/tasks/abc} e qualquer outro identificador, o nome da pasta vai entre colchetes: \texttt{app/api/tasks/[id]/}. É um \textbf{segmento dinâmico}, e o valor que veio na URL chega à função no segundo parâmetro, o \textbf{contexto}:

\begin{codigo}
\begin{Verbatim}[breaklines=true, breaksymbolleft={}, fontsize=\footnotesize]
// app/api/tasks/[id]/route.ts -- uma tarefa
import type { NextRequest } from "next/server";
import { db } from "@/lib/db";
import { usuarioAtual } from "@/lib/usuario-atual";

// GET /api/tasks/:id
export async function GET(
  request: NextRequest,
  contexto: RouteContext<"/api/tasks/[id]">,
) {
  const { id } = await contexto.params;
  const usuario = await usuarioAtual();

  const tarefa = await db.task.findFirst({
    where: { id, userId: usuario.id },
  });

  if (!tarefa) {
    return Response.json(
      { erro: "Tarefa não encontrada." },
      { status: 404 },
    );
  }
  return Response.json(tarefa);
}
\end{Verbatim}
\end{codigo}

O \texttt{RouteContext<"/api/tasks/[id]">} é um tipo que o próprio Next.js gera a partir das pastas do projeto, ao rodar \texttt{npm run dev} ou \texttt{npm run build}: ele sabe que essa rota tem um parâmetro \texttt{id}. Os parâmetros chegam como uma \textit{Promise}, por isso o \texttt{await contexto.params}; até a versão 14 do Next.js, eles eram um objeto comum, e você vai encontrar exemplos antigos sem o \texttt{await}.

A consulta é a parte mais importante. Buscar só pelo \texttt{id}, com \texttt{findUnique(\{ where: \{ id \} \})}, devolveria a tarefa de qualquer pessoa que soubesse o identificador. O \texttt{findFirst} com \texttt{id} \textbf{e} \texttt{userId} só encontra a tarefa se ela pertencer a quem pediu. Se for de outra pessoa, o resultado é o mesmo de uma tarefa que não existe: \texttt{404}. Como combinado no Capítulo 2, isso não revela a existência de tarefas alheias. A exclusão segue a mesma regra:

\begin{codigo}
\begin{Verbatim}[breaklines=true, breaksymbolleft={}, fontsize=\footnotesize]
// DELETE /api/tasks/:id
export async function DELETE(
  request: NextRequest,
  contexto: RouteContext<"/api/tasks/[id]">,
) {
  const { id } = await contexto.params;
  const usuario = await usuarioAtual();

  const { count } = await db.task.deleteMany({
    where: { id, userId: usuario.id },
  });

  if (count === 0) {
    return Response.json(
      { erro: "Tarefa não encontrada." },
      { status: 404 },
    );
  }
  return new Response(null, { status: 204 });
}
\end{Verbatim}
\end{codigo}

O \texttt{deleteMany} aceita qualquer filtro e devolve quantos registros apagou. Com o filtro por \texttt{id} e dono, o resultado é 1 (apagou) ou 0 (não havia tarefa daquele dono com aquele \texttt{id}). Fazer tudo em uma única consulta evita buscar primeiro e apagar depois. O parâmetro \texttt{request} não é usado em nenhuma das duas funções, mas precisa estar lá, porque o contexto é sempre o segundo parâmetro. A edição (\texttt{PATCH}) entra no Capítulo 7, junto com a validação de todos os campos.

\subsection{4.7 O que o Next.js faz sozinho}

Com os dois arquivos prontos, alguns comportamentos aparecem sem nenhuma linha de código a mais. Estes resultados foram obtidos com o servidor rodando:

\begin{center}
\begin{tabular}{@{}P{0.4\linewidth}P{0.5\linewidth}@{}}
\toprule
\rowcolor{headergray}
\tblheader{Requisição} & \tblheader{Resposta} \\
\midrule
\texttt{HEAD /api/tasks} & \texttt{200}, só com os cabeçalhos: o Next.js usa a função \texttt{GET} \\
\rowcolor{rowgray} \texttt{OPTIONS /api/tasks} & \texttt{204} com o cabeçalho \texttt{allow: GET, HEAD, OPTIONS, POST} \\
\texttt{PATCH /api/tasks/42} & \texttt{405 Method Not Allowed}: o arquivo não exporta \texttt{PATCH} \\
\rowcolor{rowgray} \texttt{GET /api/tarefas} & \texttt{404}: não há pasta com esse nome \\
\bottomrule
\end{tabular}
\end{center}

Ou seja: o \texttt{405} aparece quando a URL existe mas o método não, e o \texttt{404}, quando nem a URL existe, exatamente como manda o HTTP.

\subsection{4.8 Outras convenções de pastas}

Além do segmento dinâmico, o App Router tem outras convenções de nome de pasta que aparecem em APIs maiores:

\begin{center}
\begin{tabular}{@{}P{0.25\linewidth}P{0.65\linewidth}@{}}
\toprule
\rowcolor{headergray}
\tblheader{Pasta} & \tblheader{Efeito} \\
\midrule
\texttt{[id]} & um segmento variável: \texttt{/tasks/42} \\
\rowcolor{rowgray} \texttt{[...caminho]} & um ou mais segmentos: \texttt{/arquivos/a/b/c}; o parâmetro vira a lista \texttt{["a", "b", "c"]} \\
\texttt{[[...caminho]]} & como o anterior, mas também atende à URL sem nenhum segmento: \texttt{/arquivos} \\
\rowcolor{rowgray} \texttt{(grupo)} & organiza as pastas sem entrar na URL: \texttt{app/(admin)/api/users} atende \texttt{/api/users} \\
\texttt{\_privada} & ignorada pelo roteamento; serve para guardar arquivos auxiliares ao lado das rotas \\
\bottomrule
\end{tabular}
\end{center}

Quando uma URL combina com mais de uma pasta, vence a mais específica: com \texttt{app/api/tasks/[id]} e \texttt{app/api/tasks/resumo}, a URL \texttt{/api/tasks/resumo} vai para a pasta \texttt{resumo}, e não para o \texttt{[id]}.

\subsection{4.9 Testando as rotas}

O navegador só faz \texttt{GET} pela barra de endereços. Para os outros métodos, há três caminhos comuns. O primeiro é o \texttt{curl}, no terminal:

\begin{codigo}
\begin{Verbatim}[breaklines=true, breaksymbolleft={}, fontsize=\footnotesize]
$ curl http://localhost:3000/api/tasks?status=DONE
$ curl -i -X POST http://localhost:3000/api/tasks \
    -H "Content-Type: application/json" \
    -d '{"title": "Testar as rotas"}'
$ curl -i -X DELETE http://localhost:3000/api/tasks/cm3x9k2a0001
\end{Verbatim}
\end{codigo}

A opção \texttt{-X} escolhe o método, \texttt{-H} acrescenta um cabeçalho e \texttt{-d} envia o corpo. A barra invertida no fim da linha continua o comando na linha seguinte; ela funciona no terminal Git Bash, no macOS e no Linux. No PowerShell, as aspas dentro do JSON exigem cuidados extras, por isso, no Windows, prefira o segundo caminho.

O segundo é a extensão \textbf{REST Client}, do VS Code. Crie um arquivo com a extensão \texttt{.http}, escreva as requisições separadas por \texttt{\#\#\#} e clique em \textit{Send Request} acima de cada uma:

\begin{codigo}
\begin{Verbatim}[breaklines=true, breaksymbolleft={}, fontsize=\footnotesize]
### Listar as tarefas pendentes
GET http://localhost:3000/api/tasks?status=PENDING

### Criar uma tarefa
POST http://localhost:3000/api/tasks
Content-Type: application/json

{"title": "Testar as rotas"}

### Criar sem título (espera 400)
POST http://localhost:3000/api/tasks
Content-Type: application/json

{}
\end{Verbatim}
\end{codigo}

O arquivo fica no projeto e documenta os testes: qualquer pessoa da equipe consegue repeti-los. O terceiro caminho são programas com interface gráfica, como o \textbf{Postman}, assunto do Capítulo 10. Seja qual for a ferramenta, teste sempre os casos de erro, e não só o caminho feliz: título vazio, JSON quebrado, \texttt{status} inválido, identificador inexistente e exclusão repetida.

\subsection{4.10 Como está o projeto}

\begin{codigo}
\begin{Verbatim}[breaklines=true, breaksymbolleft={}, fontsize=\footnotesize]
app/
  api/
    health/route.ts     GET
    tasks/
      route.ts          GET (com ?status), POST
      [id]/route.ts     GET, DELETE
lib/
  db.ts
  usuario-atual.ts      provisório, até o Capítulo 9
testes.http             requisições para o REST Client
\end{Verbatim}
\end{codigo}

Repare que as respostas ainda trazem todos os campos da tarefa, inclusive o \texttt{userId}, que não interessa ao cliente, e que cada rota escreve a sua própria mensagem de erro. O Capítulo 5 cuida do formato dos dados que entram e saem, e o Capítulo 8 padroniza os erros.

\subsection{Erros Comuns}

\begin{itemize}
  \item \textbf{Nome da função em minúsculas.} \texttt{export async function get()} não é reconhecido, e a rota responde \texttt{405}. O nome é o do método, em maiúsculas.
  \item \textbf{\texttt{route.ts} e \texttt{page.tsx} na mesma pasta.} O Next.js acusa conflito. Rotas de API ficam em pastas próprias, como \texttt{app/api/}.
  \item \textbf{Esquecer o \texttt{await} em \texttt{contexto.params}.} Os parâmetros são uma \textit{Promise}; sem o \texttt{await}, \texttt{id} fica \texttt{undefined}.
  \item \textbf{Buscar só pelo \texttt{id}.} Sem o \texttt{userId} no filtro, qualquer pessoa lê ou apaga a tarefa de outra. Essa falha é tão comum que lidera a lista de riscos de APIs da OWASP, estudada no Capítulo 12.
  \item \textbf{Deixar o JSON inválido virar 500.} Trate o erro de \texttt{request.json()} e responda \texttt{400}.
  \item \textbf{Tratar a \textit{query string} como número.} \texttt{searchParams.get("page")} devolve o texto \texttt{"2"}; converta com \texttt{Number()} e confira o resultado.
  \item \textbf{Mandar corpo com o status 204.} O 204 significa ``sem conteúdo''; use \texttt{new Response(null, \{ status: 204 \})}.
  \item \textbf{Esquecer o \texttt{return}.} Se a função monta a resposta de erro mas não a devolve, o código continua e executa a consulta mesmo assim.
\end{itemize}
